\documentclass[12pt,a4paper]{article}
\usepackage[utf8]{inputenc}
\usepackage[brazil]{babel}
\usepackage{geometry}
\usepackage{listings}
\usepackage{xcolor}
\usepackage{tikz}
\usepackage{amsmath}
\usepackage{hyperref}
\usepackage{fancyhdr}

\geometry{margin=2.5cm}
\setlength{\headheight}{14.5pt}

% Configuracao de sintaxe Pascal (comentarios estilo (* *))
\lstdefinestyle{pascal}{
    language=Pascal,
    basicstyle=\ttfamily\footnotesize,
    keywordstyle=\color{blue}\bfseries,
    commentstyle=\color{green!60!black}\itshape,
    stringstyle=\color{red},
    morecomment=[s]{(*}{*)},
    numbers=left,
    numberstyle=\tiny\color{gray},
    stepnumber=5,
    numbersep=8pt,
    showstringspaces=false,
    breaklines=true,
    frame=single,
    backgroundcolor=\color{gray!5},
    xleftmargin=15pt,
    framexleftmargin=15pt
}

\usetikzlibrary{automata,positioning,arrows.meta,shapes}

\title{\textbf{Conversor AFN para AFD}\\
\Large Estruturas de Dados e Algoritmo Principal\\
\large Documentacao Tecnica - Henrique}
\author{Henrique Freitas}
\date{Dezembro 2025}

\pagestyle{fancy}
\fancyhf{}
\fancyhead[L]{Conversor AFN para AFD}
\fancyhead[R]{Henrique - Estruturas e Conversao}
\fancyfoot[C]{\thepage}

\begin{document}

\maketitle
\tableofcontents
\newpage

\section{Introducao}

Este documento detalha as secoes do codigo \texttt{MainForm.pas} sob responsabilidade de \textbf{Henrique}:

\begin{itemize}
    \item \textbf{Parte 1:} Linhas 1-2254 (Estruturas de Dados, Interface, Parsing)
    \item \textbf{Parte 2:} Linhas 2959-3795 (Algoritmo AFN para AFD)
\end{itemize}

\subsection{Visao Geral do Projeto}

O projeto implementa um conversor completo de automatos finitos:
\begin{enumerate}
    \item \textbf{AFN-$\varepsilon$} (com epsilon) para \textbf{AFN} (sem epsilon)
    \item \textbf{AFN} (nao-deterministico) para \textbf{AFD} (deterministico)
    \item \textbf{AFD} para \textbf{AFD Minimizado}
\end{enumerate}

\textbf{Tecnologia:} Free Pascal + Lazarus (GUI nativa)

\newpage
\section{Parte 1: Estruturas de Dados e Interface (L1-2254)}

\subsection{Visao Geral da Parte 1}

A Parte 1 contem:
\begin{itemize}
    \item Tipos fundamentais (TTransition, TTransitionArray)
    \item Classe TStringSet (L299-475)
    \item Variaveis globais (L135-150)
    \item FormCreate - Inicializacao (L488-571)
    \item LoadTestFiles - Scan de arquivos (L586-643)
    \item OnTestFileSelected - Handler ComboBox (L645-679)
    \item ConvertJSONToText - Parse arquivo JSON (L693-812)
    \item ConvertJSONTextToText - Parse string JSON (L829-918)
    \item LoadSampleFile - Exemplo padrao (L952-1043)
    \item btnLoadFileClick - Carregar arquivo (L1085-1170)
    \item btnClearClick - Limpar interface (L1237-1360)
    \item btnRemoveEpsilonClick - Handler epsilon (L1395-1566)
    \item btnUseAsInputClick - Copiar resultado (L1603-1720)
    \item btnConvertClick - Handler conversao (L1863-2068)
\end{itemize}

\subsection{Tipos Fundamentais (L39-50)}

\subsubsection{TTransition - Record de Transicao (L39-43)}

\textbf{Referencia:} \texttt{MainForm.pas} linhas 39-43

\begin{lstlisting}[style=pascal, caption={Definicao de TTransition}]
TTransition = record
  FromState: string;  (* Estado de origem *)
  Symbol: string;     (* Simbolo da transicao *)
  ToState: string;    (* Estado de destino *)
end;
\end{lstlisting}

\textbf{Explicacao:}
\begin{itemize}
    \item Representa uma transicao: $\delta(q_i, a) = q_j$
    \item \texttt{FromState} = $q_i$, \texttt{Symbol} = $a$, \texttt{ToState} = $q_j$
\end{itemize}

\subsubsection{TTransitionArray - Array Dinamico (L49)}

\begin{lstlisting}[style=pascal]
TTransitionArray = array of TTransition;
\end{lstlisting}

Array dinamico que armazena todas as transicoes do automato.

\subsection{Classe TFormMain - Formulario Principal (L50-134)}

\textbf{Referencia:} \texttt{MainForm.pas} linhas 50-134

\subsubsection{Declaracao da Classe (L76-130)}

A classe \texttt{TFormMain} herda de \texttt{TForm} e gerencia toda a interface grafica:

\textbf{Componentes Visuais:}
\begin{itemize}
    \item \textbf{Botoes:} btnConvert, btnLoadFile, btnClear, btnMinimize, btnRemoveEpsilon, btnUseAsInput
    \item \textbf{Editores de texto:} memoInput, memoOutput, memoNFAOutput, memoMinOutput
    \item \textbf{Paineis de desenho:} PaintBoxEpsilonNFA, PaintBoxNFA, PaintBoxDFA, PaintBoxMinDFA
    \item \textbf{Abas:} TabEpsilonNFA, TabNFA, TabDFA, TabMinDFA
    \item \textbf{Dialogo:} OpenDialog (para selecionar arquivos)
\end{itemize}

\textbf{Event Handlers:}
\begin{itemize}
    \item btnClearClick - Limpar interface
    \item btnConvertClick - Converter AFN para AFD
    \item btnLoadFileClick - Carregar arquivo
    \item btnMinimizeClick - Minimizar AFD
    \item btnRemoveEpsilonClick - Remover epsilon-transicoes
    \item btnUseAsInputClick - Copiar resultado para entrada
    \item FormCreate - Inicializacao do formulario
    \item PaintBoxXXXPaint - Desenhar diagramas
\end{itemize}

\subsection{Variaveis Globais dos Automatos (L135-165)}

\textbf{Referencia:} \texttt{MainForm.pas} linhas 135-165

\begin{lstlisting}[style=pascal, caption={Variaveis Globais}]
(* AFN-epsilon: entrada original *)
EpsilonNFAAlphabet, EpsilonNFAStates: TStringList;
EpsilonNFAInitials, EpsilonNFAFinals: TStringList;
EpsilonNFATransitions: TTransitionArray;

(* AFN: sem epsilon *)
NFAAlphabet, NFAStates, NFAInitials, NFAFinals: TStringList;
NFATransitions: TTransitionArray;

(* AFD: deterministico *)
DFAStates, DFAFinals: TStringList;
DFAInitial: string;  (* AFD tem apenas UM inicial *)
DFATransitions: TTransitionArray;

(* AFD Minimizado *)
MinDFAStates, MinDFAFinals: TStringList;
MinDFAInitial: string;
MinDFATransitions: TTransitionArray;
\end{lstlisting}

\subsection{Declaracoes de Metodos Privados (L166-298)}

\textbf{Referencia:} \texttt{MainForm.pas} linhas 166-298

A secao privada da classe declara os metodos principais de conversao:

\subsubsection{RemoveEpsilonTransitions (L174)}

\textbf{Proposito:} Remove epsilon-transicoes de um AFN-$\varepsilon$

\textbf{Algoritmo:}
\begin{enumerate}
    \item Calcula epsilon-fecho ($\varepsilon$-closure) de cada estado
    \item Para cada estado e simbolo, adiciona transicoes para todos destinos alcancaveis
    \item Ajusta estados finais (estado e final se seu $\varepsilon$-fecho contem algum final)
\end{enumerate}

\textbf{Complexidade:} $O(n^2 \cdot m)$ onde $n$ = estados, $m$ = transicoes

\subsubsection{ConvertAFNtoAFD (L190)}

\textbf{Proposito:} Converte AFN em AFD usando construcao de subconjuntos

\textbf{Algoritmo (Subset Construction):}
\begin{enumerate}
    \item Estado inicial do AFD = conjunto dos estados iniciais do AFN
    \item Para cada estado do AFD (conjunto de estados do AFN):
    \begin{itemize}
        \item Para cada simbolo do alfabeto:
        \item Calcula conjunto de estados alcancaveis
        \item Cria novo estado do AFD se necessario
        \item Adiciona transicao
    \end{itemize}
    \item Estado do AFD e final se contem algum estado final do AFN
\end{enumerate}

\textbf{Complexidade:} $O(2^n \cdot |\Sigma|)$ no pior caso (explosao exponencial)

\subsubsection{MinimizeDFA (L207)}

\textbf{Proposito:} Minimiza um AFD usando algoritmo de Myhill-Nerode

\textbf{Algoritmo (Particao-Refinamento):}
\begin{enumerate}
    \item Particao inicial: (finais, nao-finais)
    \item Repetir ate nao haver mudancas:
    \begin{itemize}
        \item Para cada particao P e simbolo a:
        \item Se estados de P vao para particoes diferentes com 'a', divide P
    \end{itemize}
    \item Cada particao final vira um estado do MinDFA
\end{enumerate}

\textbf{Complexidade:} $O(n \cdot m \cdot \log n)$ onde $n$ = estados, $m$ = transicoes

\subsubsection{Metodos Auxiliares (L222-298)}

\textbf{LoadSampleFile:} Carrega arquivo de exemplo \texttt{sample\_afn.txt}

\textbf{LoadTestFiles:} Popula ComboBox com arquivos do diretorio \texttt{testes/}

\textbf{OnTestFileSelected:} Handler para selecao de arquivo de teste

\textbf{ConvertJSONToText:} Converte arquivo JSON para formato TXT

\textbf{ConvertJSONTextToText:} Converte string JSON para formato TXT

\textbf{DrawAutomaton:} Desenha diagrama de automato em TCanvas
\begin{itemize}
    \item Layout: 4 estados por linha
    \item Espacamento: 100px horizontal, 80px vertical
    \item Estados: circulos de raio 20px
    \item Estados finais: circulo duplo com fundo amarelo
    \item Transicoes: linhas com setas e rotulos
\end{itemize}

\subsection{TStringSet - Classe de Conjuntos (L299-475)}

\textbf{Referencia:} \texttt{MainForm.pas} linhas 299-475

\subsubsection{Declaracao da Classe (L299-311)}

\begin{lstlisting}[style=pascal, caption={Classe TStringSet}]
TStringSet = class
  private
    FList: TStringList;
  public
    constructor Create;
    destructor Destroy; override;
    procedure Add(const s: string);
    function Contains(const s: string): Boolean;
    function ToString: string; override;
    function Clone: TStringSet;
    function IsEmpty: Boolean;
    function Count: Integer;
    function Item(i: Integer): string;
end;
\end{lstlisting}

\textbf{Por Que TStringSet?}

Construcao de subconjuntos gera muitos conjuntos de estados. Precisamos evitar duplicatas. TStringList com \texttt{Duplicates := dupIgnore} resolve!

\subsubsection{Implementacao dos Metodos (L333-475)}

\textbf{Create (L333-337):} Inicializa com Sorted=True e Duplicates=dupIgnore

\textbf{Destroy (L339-343):} Libera FList

\textbf{Add (L360-367):} Adiciona elemento (ignora duplicatas automaticamente)

\textbf{Contains (L377-381):} Verifica se elemento existe (O(log n) com busca binaria)

\textbf{Clone (L398-410):} Cria copia independente do conjunto

\textbf{ToString (L431-445):} Retorna formato "\{q0,q1,q2\}"

\textbf{IsEmpty (L414-418):} Verifica se conjunto vazio

\textbf{Count (L449-453):} Retorna numero de elementos

\textbf{Item (L464-475):} Acessa elemento por indice

\subsection{Inicializacao da Interface (L488-571)}

\textbf{Referencia:} \texttt{MainForm.pas} linhas 488-571

\begin{lstlisting}[style=pascal, caption={FormCreate}]
procedure TFormMain.FormCreate(Sender: TObject);
begin
  (* Criar listas para alfabeto, estados, etc. *)
  NFAAlphabet := TStringList.Create;
  NFAStates := TStringList.Create;
  NFAInitials := TStringList.Create;
  NFAFinals := TStringList.Create;
  DFAStates := TStringList.Create;
  DFAFinals := TStringList.Create;
  MinDFAStates := TStringList.Create;
  MinDFAFinals := TStringList.Create;
  
  (* Carregar arquivos de teste *)
  LoadTestFiles;
  LoadSampleFile;
end;
\end{lstlisting}

\subsection{Parser JSON para TXT (L693-918)}

O sistema aceita dois formatos: TXT e JSON. Duas funcoes realizam a conversao:

\begin{itemize}
    \item \textbf{ConvertJSONToText} (L693-812): Le arquivo JSON do disco
    \item \textbf{ConvertJSONTextToText} (L829-918): Converte string JSON direta
\end{itemize}

\subsection{LoadSampleFile - Exemplo Padrao (L952-1043)}

\textbf{Referencia:} \texttt{MainForm.pas} linhas 952-1043

Carrega arquivo sample\_afn.txt se existir, senao cria exemplo hardcoded. O exemplo padrao e um AFN que aceita strings terminadas em "ab".

\subsection{btnLoadFileClick - Carregar Arquivo (L1085-1170)}

\textbf{Referencia:} \texttt{MainForm.pas} linhas 1085-1170

Event handler do botao "Carregar Arquivo". Abre dialogo nativo (TOpenDialog) para selecionar arquivo .txt ou .json do sistema de arquivos.

\begin{lstlisting}[style=pascal, caption={btnLoadFileClick}]
procedure TFormMain.btnLoadFileClick(Sender: TObject);
begin
  OpenDialog.Filter := 'Arquivos de automato (*.txt;*.json)|...';
  OpenDialog.DefaultExt := 'txt';
  
  if OpenDialog.Execute then
  begin
    try
      memoInput.Lines.LoadFromFile(OpenDialog.FileName);
      edtFilePath.Text := OpenDialog.FileName;
    except
      on E: Exception do
        ShowMessage('Erro ao carregar arquivo: ' + E.Message);
    end;
  end;
end;
\end{lstlisting}

\subsection{btnClearClick - Limpar Interface (L1237-1360)}

\textbf{Referencia:} \texttt{MainForm.pas} linhas 1237-1360

Limpa completamente a interface e todos os dados, restaurando ao estado inicial.

\textbf{Operacoes:}
\begin{itemize}
    \item Limpa todos os memos (Input, Output, NFAOutput, MinOutput)
    \item Limpa todas as TStringList
    \item Libera arrays de transicoes (SetLength 0)
    \item Desabilita botoes dependentes
    \item Redesenha diagramas vazios
\end{itemize}

\subsection{btnRemoveEpsilonClick - Handler Epsilon (L1395-1566)}

\textbf{Referencia:} \texttt{MainForm.pas} linhas 1395-1566

Event handler do botao "AFN-epsilon para AFN". Detecta se entrada e JSON e converte temporariamente. Verifica se ha epsilon-transicoes antes de processar. Invoca RemoveEpsilonTransitions se validacao passar.

\subsection{btnUseAsInputClick - Copiar Resultado (L1603-1720)}

\textbf{Referencia:} \texttt{MainForm.pas} linhas 1603-1720

Copia AFN resultante (de memoNFAOutput) para memoInput, permitindo workflow sequencial.

\begin{lstlisting}[style=pascal, caption={btnUseAsInputClick}]
procedure TFormMain.btnUseAsInputClick(Sender: TObject);
var i: Integer;
begin
  memoInput.Lines.Clear;
  memoInput.Lines.Add(NFAAlphabet.CommaText.Replace(',', ' '));
  memoInput.Lines.Add(NFAStates.CommaText.Replace(',', ' '));
  memoInput.Lines.Add(NFAInitials.CommaText.Replace(',', ' '));
  memoInput.Lines.Add(NFAFinals.CommaText.Replace(',', ' '));
  
  for i := 0 to High(NFATransitions) do
    memoInput.Lines.Add(NFATransitions[i].FromState + ' ' + 
                       NFATransitions[i].Symbol + ' ' + 
                       NFATransitions[i].ToState);
  
  ShowMessage('AFN sem epsilon copiado para a entrada!');
end;
\end{lstlisting}

\subsection{btnConvertClick - Handler Conversao (L1863-2068)}

\textbf{Referencia:} \texttt{MainForm.pas} linhas 1863-2068

Event handler principal do botao "AFN para AFD". Detecta JSON e converte automaticamente. Se detectar epsilon-transicoes, executa workflow completo: AFN-epsilon para AFN para AFD. Invoca ConvertAFNtoAFD ao final.

\newpage
\section{Parte 2: Conversao AFN para AFD (L2958-3795)}

\subsection{Teoria: Construcao de Subconjuntos}

\subsubsection{A Grande Ideia}

\begin{center}
\fbox{
\begin{minipage}{0.9\textwidth}
\textbf{Principio Fundamental:}\\
Cada estado do AFD representa um \textbf{conjunto de estados} do AFN onde voce poderia estar simultaneamente.
\end{minipage}
}
\end{center}

\subsubsection{Motivacao: Por Que Precisamos Disto?}

\textbf{Problema:} AFN pode ter multiplas transicoes com mesmo simbolo.

\textbf{Exemplo:}
\begin{itemize}
    \item $\delta(q_0, a) = \{q_1, q_2\}$ -- AFN vai para q1 \textbf{E} q2 simultaneamente!
    \item AFD exige: $\delta(q_0, a) = q_X$ -- apenas UM destino
\end{itemize}

\textbf{Solucao:} Criar estado $q_X = \{q_1, q_2\}$ no AFD.

\subsubsection{Por Que TStringSet e Essencial?}

\textbf{Problema 1 - Duplicatas:}
\begin{itemize}
    \item Ao calcular $\bigcup_{q \in S} \delta(q, a)$, mesmos estados aparecem multiplas vezes
    \item TStringSet com \texttt{Duplicates := dupIgnore} elimina automaticamente!
\end{itemize}

\textbf{Problema 2 - Comparacao de Conjuntos:}
\begin{itemize}
    \item $\{q_0, q_1\}$ e $\{q_1, q_0\}$ sao o mesmo estado!
    \item TStringSet com \texttt{Sorted := True} garante ordem alfabetica
    \item ToString() sempre retorna mesma string para conjuntos equivalentes
\end{itemize}

\textbf{Problema 3 - Performance:}
\begin{itemize}
    \item Busca linear: $O(n)$
    \item Busca binaria (Sorted=True): $O(\log n)$
    \item Com milhares de estados, diferenca e gigantesca!
\end{itemize}

\subsection{Estrutura do Algoritmo ConvertAFNtoAFD (L2958-3795)}

\textbf{Referencia:} \texttt{MainForm.pas} linhas 2958-3795

\begin{lstlisting}[style=pascal, caption={Declaracoes}]
procedure TFormMain.ConvertAFNtoAFD;
var
  (* Variaveis para parsing do AFN *)
  Alphabet, States, Initials, Finals: TStringList;
  Transitions: array of TTransition;
  parts: TStringList;
  i, j, k, tStart: Integer;
  
  (* Variaveis para construcao do AFD *)
  localDFAStates: TStringList;   (* Estados do AFD *)
  workQ: TStringList;             (* Fila de trabalho BFS *)
  dfaMap: TStringList;            (* Mapa nome para TStringSet *)
  isFinal: TStringList;           (* Estados finais do AFD *)
  
  curSet, nextSet, cloneSet: TStringSet;
  sym, key: string;
  localDFATransitions: TTransitionArray;
  hasFinal, hasEpsilon: Boolean;
\end{lstlisting}

\subsection{Declaracao de Variaveis (L2959-2978)}

\textbf{Referencia:} \texttt{MainForm.pas} linhas 2959-2978

O algoritmo declara dois grupos de variaveis:

\textbf{Grupo 1 - Parsing do AFN de entrada:}
\begin{itemize}
    \item \texttt{Alphabet, States, Initials, Finals}: TStringList para armazenar componentes do AFN
    \item \texttt{Transitions}: Array dinamico de transicoes do AFN
    \item \texttt{parts}: TStringList para quebrar linhas em tokens
    \item \texttt{i, j, k, tStart}: Indices para loops
\end{itemize}

\textbf{Grupo 2 - Construcao do AFD:}
\begin{itemize}
    \item \texttt{localDFAStates}: Lista de nomes dos estados do AFD (strings como "(q0,q1)")
    \item \texttt{workQ}: Fila de trabalho BFS (estados que precisam ser processados)
    \item \texttt{dfaMap}: Mapeamento nome $\rightarrow$ conjunto real (TStringSet)
    \item \texttt{isFinal}: Lista de estados finais do AFD
    \item \texttt{curSet, nextSet, cloneSet}: TStringSet para manipular conjuntos
    \item \texttt{localDFATransitions}: Array de transicoes do AFD
\end{itemize}

\subsection{Limpeza e Inicializacao (L2980-3005)}

\textbf{Referencia:} \texttt{MainForm.pas} linhas 2980-3005

\textbf{Por que limpar?} Evitar dados misturados de conversoes anteriores.

\textbf{Procedimento:}
\begin{enumerate}
    \item Verifica se objetos existem com \texttt{Assigned()}
    \item Libera memoria com \texttt{Free()}
    \item Cria novas instancias vazias
\end{enumerate}

\textbf{Estruturas recriadas:}
\begin{itemize}
    \item NFAAlphabet, NFAStates, NFAInitials, NFAFinals (para guardar AFN)
    \item DFAStates, DFAFinals (para construir AFD)
\end{itemize}

\subsection{Validacao de Entrada (L3007-3018)}

\textbf{Referencia:} \texttt{MainForm.pas} linhas 3007-3018

Verifica se entrada tem pelo menos 4 linhas:
\begin{enumerate}
    \item Alfabeto
    \item Estados
    \item Iniciais
    \item Finais
\end{enumerate}

Se invalida: Mostra mensagem e aborta com \texttt{Exit}.

\subsection{Parse do Alfabeto (L3033-3058)}

\textbf{Referencia:} \texttt{MainForm.pas} linhas 3033-3058

\textbf{Objetivo:} Ler simbolos do alfabeto da linha 0.

\textbf{Logica:}
\begin{lstlisting}[style=pascal]
parts.DelimitedText := Trim(memoInput.Lines[0]);
for i := 0 to parts.Count - 1 do
  if (simbolo nao e epsilon) then
    Alphabet.Add(parts[i]);
\end{lstlisting}

\textbf{Filtragem de epsilon:} AFD nao pode ter $\varepsilon$-transicoes, entao filtra:
\begin{itemize}
    \item '$\varepsilon$' (Unicode)
    \item 'epsilon' (palavra)
    \item 'e' (abreviacao)
    \item '\&' (alternativa)
\end{itemize}

\subsection{Parse dos Estados (L3067-3075)}

\textbf{Referencia:} \texttt{MainForm.pas} linhas 3067-3075

\textbf{Objetivo:} Ler lista de estados da linha 1.

Cada estado e uma string (ex: "q0", "q1", "q2").

Armazena em:
\begin{itemize}
    \item \texttt{States}: Lista local
    \item \texttt{NFAStates}: Campo da classe (usado depois)
\end{itemize}

\subsection{Parse dos Estados Iniciais (L3084-3092)}

\textbf{Referencia:} \texttt{MainForm.pas} linhas 3084-3092

\textbf{Objetivo:} Ler estados iniciais da linha 2.

\textbf{Importante:} AFN pode ter multiplos iniciais!

No AFD: estado inicial = \textbf{conjunto} dos iniciais do AFN.

\subsection{Parse dos Estados Finais (L3101-3109)}

\textbf{Referencia:} \texttt{MainForm.pas} linhas 3101-3109

\textbf{Objetivo:} Ler estados finais da linha 3.

\textbf{Regra para AFD:} Estado do AFD e final $\Leftrightarrow$ contem algum estado final do AFN.

\subsection{Parse das Transicoes (L3130-3177)}

\textbf{Referencia:} \texttt{MainForm.pas} linhas 3130-3177

\textbf{Formato esperado:} \texttt{origem simbolo destino}

\textbf{Algoritmo:}
\begin{lstlisting}[style=pascal]
for i := 4 to memoInput.Lines.Count - 1 do
  parts.DelimitedText := linha;
  if parts.Count >= 3 then
    Transicao := (parts[0], parts[1], parts[2]);
\end{lstlisting}

\textbf{Deteccao de epsilon:} Se encontrar $\varepsilon$-transicao, marca flag \texttt{hasEpsilon}.

\subsection{Aviso de Epsilon (L3179-3191)}

\textbf{Referencia:} \texttt{MainForm.pas} linhas 3179-3191

Se \texttt{hasEpsilon = True}:
\begin{itemize}
    \item Mostra dialogo de aviso
    \item Orienta usuario: AFN-$\varepsilon$ $\rightarrow$ AFN $\rightarrow$ AFD
    \item Continua execucao (resultado pode estar incorreto)
\end{itemize}

\textbf{Por que avisar?} AFD nao pode ter $\varepsilon$ por definicao!

\subsubsection{Filosofia de Validacao: Avisar vs Bloquear}

\textbf{Alternativa 1 - Bloquear execucao:}
\begin{lstlisting}[style=pascal]
if hasEpsilon then
begin
  ShowMessage('ERRO: AFD nao aceita epsilon!');
  Exit;  (* Aborta conversao *)
end;
\end{lstlisting}

\textbf{Alternativa 2 - Avisar e continuar (implementado):}
\begin{lstlisting}[style=pascal]
if hasEpsilon then
  ShowMessage('AVISO: Entrada tem epsilon. Use RemoveEpsilon primeiro.');
(* Continua execucao *)
\end{lstlisting}

\textbf{Por que escolhemos Alternativa 2?}
\begin{enumerate}
    \item \textbf{Flexibilidade:} Usuario pode querer testar mesmo assim
    \item \textbf{Debugging:} Permite ver resultado incorreto para comparacao
    \item \textbf{Workflow educacional:} Mostra o que acontece com epsilon
    \item \textbf{Nao-destrutivo:} Nao perde trabalho do usuario
\end{enumerate}

\textbf{Trade-off:}
\begin{itemize}
    \item \textbf{Pro:} Usuario tem controle e visibilidade
    \item \textbf{Contra:} Usuario pode ignorar aviso e obter resultado errado
\end{itemize}

\textbf{Decisao de design:} Em ferramenta educacional, preferimos avisar e ensinar ao inves de bloquear.

\subsection{Logging da Entrada (L3206-3217)}

\textbf{Referencia:} \texttt{MainForm.pas} linhas 3206-3217

Imprime no console (WriteLn) o AFN de entrada para debugging:
\begin{itemize}
    \item Alfabeto
    \item Estados
    \item Iniciais
    \item Finais
    \item Lista de transicoes
\end{itemize}

\subsection{Criacao das Estruturas BFS (L3251-3257)}

\textbf{Referencia:} \texttt{MainForm.pas} linhas 3251-3257

\textbf{Estruturas criadas:}

\begin{tabular}{ll}
\texttt{localDFAStates} & Lista de estados do AFD (nomes) \\
\texttt{workQ} & Fila BFS (estados a processar) \\
\texttt{dfaMap} & Mapa nome $\rightarrow$ TStringSet \\
\texttt{isFinal} & Estados finais do AFD \\
\end{tabular}

\textbf{Importante:} \texttt{dfaMap.OwnsObjects := True}

$\rightarrow$ libera TStringSets automaticamente!

\subsubsection{Por Que BFS e Nao DFS?}

\textbf{Razao 1 - Ordem Natural:}
\begin{itemize}
    \item BFS processa estados por "distancia" do inicial
    \item Resultado mais intuitivo para visualizacao
    \item Estados principais aparecem primeiro no output
\end{itemize}

\textbf{Razao 2 - Memoria Controlada:}
\begin{itemize}
    \item DFS: Pilha pode crescer exponencialmente em automatos complexos
    \item BFS: Fila cresce de forma mais previsivel
    \item Importante para automatos com milhares de estados
\end{itemize}

\textbf{Razao 3 - Deteccao de Loops:}
\begin{itemize}
    \item BFS detecta estados repetidos naturalmente
    \item DFS pode ficar preso em ciclos se nao implementado corretamente
\end{itemize}

\subsubsection{Por Que OwnsObjects = True?}

\textbf{Problema sem OwnsObjects:}
\begin{lstlisting}[style=pascal]
dfaMap.AddObject('\{q0,q1\}', set1);
dfaMap.AddObject('\{q2,q3\}', set2);
dfaMap.Free;  (* Libera dfaMap mas NAO os TStringSets! *)
(* MEMORY LEAK! set1 e set2 ainda estao na memoria! *)
\end{lstlisting}

\textbf{Solucao com OwnsObjects:}
\begin{lstlisting}[style=pascal]
dfaMap.OwnsObjects := True;
dfaMap.AddObject('\{q0,q1\}', set1);
dfaMap.Free;  (* Libera dfaMap E todos os TStringSets! *)
(* Sem memory leak! *)
\end{lstlisting}

\textbf{Por que isso importa?}
\begin{itemize}
    \item AFD com 100 estados = 100 TStringSet para liberar manualmente!
    \item Esqueceu um Free()? Memory leak.
    \item OwnsObjects = True $\rightarrow$ gerenciamento automatico!
\end{itemize}

\subsection{Estado Inicial do AFD (L3273-3283)}

\textbf{Referencia:} \texttt{MainForm.pas} linhas 3273-3283

\textbf{Regra:} Estado inicial do AFD = \textbf{conjunto} dos iniciais do AFN.

\textbf{Algoritmo:}
\begin{lstlisting}[style=pascal]
curSet := TStringSet.Create;
for i := 0 to Initials.Count - 1 do
  curSet.Add(Initials[i]);
  
localDFAStates.Add(GetSetName(curSet));
workQ.Add(GetSetName(curSet));
dfaMap.AddObject(GetSetName(curSet), curSet);
\end{lstlisting}

\textbf{Exemplo:} Se AFN tem iniciais (q0, q1) $\rightarrow$ estado inicial do AFD = "\{q0,q1\}"

\subsection{Loop Principal BFS (L3296-3540)}

\textbf{Referencia:} \texttt{MainForm.pas} linhas 3296-3540

\textbf{Algoritmo de Busca em Largura:}

\begin{lstlisting}[style=pascal]
while workQ.Count > 0 do
  (* 1. Remover estado da fila *)
  curSet := ProximoEstado();
  
  (* 2. Verificar se e final *)
  if curSet contem algum final do AFN then
    Marcar como final;
  
  (* 3. Para cada simbolo do alfabeto *)
  for cada simbolo do
    nextSet := CalcularDestinos(curSet, simbolo);
    CriarTransicao(curSet, simbolo, nextSet);
    
    (* 4. Se nextSet e novo, adicionar a fila *)
    if nextSet e novo then
      workQ.Add(nextSet);
end;
\end{lstlisting}

\textbf{Invariante:} Todos estados em \texttt{localDFAStates} foram ou serao processados.

\textbf{Terminacao:} Numero de estados $\leq 2^n$ (finito).

\subsubsection{Exemplo Pratico Passo-a-Passo}

\textbf{AFN de entrada:}
\begin{itemize}
    \item Estados: $\{q_0, q_1, q_2\}$
    \item Alfabeto: $\{a, b\}$
    \item Iniciais: $\{q_0\}$
    \item Finais: $\{q_2\}$
    \item Transicoes: $(q_0, a, q_0)$, $(q_0, a, q_1)$, $(q_1, b, q_2)$
\end{itemize}

\textbf{Execucao do BFS:}

\textbf{Iteracao 1:}
\begin{enumerate}
    \item \texttt{workQ = ["\{q0\}"]}, remove "\{q0\}"
    \item Nao e final (q0 nao e final)
    \item Com 'a': destinos = \{q0, q1\} $\rightarrow$ transicao (\{q0\}, a, \{q0,q1\})
    \item Com 'b': destinos = \{\} $\rightarrow$ transicao (\{q0\}, b, \{\})
    \item Novo estado descoberto: "\{q0,q1\}" $\rightarrow$ adiciona a fila
    \item \texttt{workQ = ["\{q0,q1\}"]}
\end{enumerate}

\textbf{Iteracao 2:}
\begin{enumerate}
    \item \texttt{workQ = ["\{q0,q1\}"]}, remove "\{q0,q1\}"
    \item Nao e final (nem q0 nem q1 sao finais)
    \item Com 'a': destinos = \{q0, q1\} $\rightarrow$ transicao (\{q0,q1\}, a, \{q0,q1\})
    \item Com 'b': destinos = \{q2\} $\rightarrow$ transicao (\{q0,q1\}, b, \{q2\})
    \item Novo estado descoberto: "\{q2\}" $\rightarrow$ adiciona a fila
    \item \texttt{workQ = ["\{q2\}"]}
\end{enumerate}

\textbf{Iteracao 3:}
\begin{enumerate}
    \item \texttt{workQ = ["\{q2\}"]}, remove "\{q2\}"
    \item E final! (q2 e final)
    \item Com 'a': destinos = \{\} $\rightarrow$ transicao (\{q2\}, a, \{\})
    \item Com 'b': destinos = \{\} $\rightarrow$ transicao (\{q2\}, b, \{\})
    \item Nenhum novo estado
    \item \texttt{workQ = []}
\end{enumerate}

\textbf{Fim!} AFD resultante tem 3 estados: \{q0\}, \{q0,q1\}, \{q2\}

\subsection{Remocao da Fila (L3308-3311)}

\textbf{Referencia:} \texttt{MainForm.pas} linhas 3308-3311

\begin{lstlisting}[style=pascal]
key := workQ[0];
workQ.Delete(0);
curSet := TStringSet(dfaMap.Objects[...]);
\end{lstlisting}

\textbf{FIFO:} Remove do inicio (BFS), nao do fim.

\subsection{Verificacao de Estado Final (L3326-3330)}

\textbf{Referencia:} \texttt{MainForm.pas} linhas 3326-3330

\textbf{Regra:} Estado AFD e final $\Leftrightarrow$ \texttt{curSet} contem algum final do AFN.

\begin{lstlisting}[style=pascal]
hasFinal := False;
for i := 0 to Finals.Count - 1 do
  if curSet.Contains(Finals[i]) then 
    hasFinal := True;
if hasFinal then isFinal.Add(key);
\end{lstlisting}

\textbf{Justificativa:} Se caminho leva a estado final do AFN, o conjunto todo e final.

\subsection{Loop de Simbolos (L3351-3525)}

\textbf{Referencia:} \texttt{MainForm.pas} linhas 3351-3525

Para cada simbolo do alfabeto, calcula proximos estados:

\textbf{Algoritmo:}
\begin{lstlisting}[style=pascal]
for i := 0 to Alphabet.Count - 1 do
  sym := Alphabet[i];
  nextSet := TStringSet.Create;
  
  (* Uniao dos destinos *)
  for cada q em curSet do
    for cada transicao (q, sym, r) do
      nextSet.Add(r);
  
  (* Criar transicao AFD *)
  Transicao := (curSet, sym, nextSet);
\end{lstlisting}

\textbf{Formula:} $\delta_{AFD}(S, a) = \bigcup_{q \in S} \delta_{AFN}(q, a)$

\subsection{Calculo de Destinos (L3356-3373)}

\textbf{Referencia:} \texttt{MainForm.pas} linhas 3356-3373

\begin{lstlisting}[style=pascal]
for j := 0 to curSet.Count - 1 do
  for k := 0 to High(Transitions) do
    if (Transitions[k].FromState = curSet.Item(j))
       and (Transitions[k].Symbol = sym) then
      nextSet.Add(Transitions[k].ToState);
\end{lstlisting}

\textbf{Complexidade:} $O(|curSet| \cdot m)$ onde $m$ = transicoes do AFN.

\subsection{Tratamento de Conjunto Vazio (L3375-3379)}

\textbf{Referencia:} \texttt{MainForm.pas} linhas 3375-3379

Se \texttt{nextSet.IsEmpty}:
\begin{itemize}
    \item AFD vai para estado de erro "\{\}"
    \item Estado de erro NAO e adicionado ao AFD (dead state implicito)
\end{itemize}

\textbf{Exemplo:} (q0) com 'b' e nenhuma transicao existe $\rightarrow$ "\{\}"

\subsubsection{Por Que Nao Adicionar Estado de Erro?}

\textbf{Razao 1 - Tamanho do AFD:}
\begin{itemize}
    \item Estado de erro tem transicoes para si mesmo com TODOS os simbolos
    \item Em alfabeto grande, adiciona muitas transicoes inuteis
    \item Exemplo: alfabeto de 26 letras = 26 transicoes extras!
\end{itemize}

\textbf{Razao 2 - Visualizacao:}
\begin{itemize}
    \item Estado de erro polui diagrama
    \item Usuario sabe: "sem transicao = rejeita"
    \item Mais limpo omitir
\end{itemize}

\textbf{Razao 3 - Minimizacao:}
\begin{itemize}
    \item Algoritmo de minimizacao remove estado de erro de qualquer forma
    \item Por que adicionar so para remover depois?
\end{itemize}

\textbf{Trade-off:}
\begin{itemize}
    \item \textbf{Com estado erro:} AFD teoricamente completo
    \item \textbf{Sem estado erro:} AFD parcial (mas funcionalmente identico)
\end{itemize}

\textbf{Decisao de implementacao:} Omitir estado de erro para AFD mais limpo e legivel.

\subsection{Registro de Transicao (L3395-3409)}

\textbf{Referencia:} \texttt{MainForm.pas} linhas 3395-3409

Expande array dinamico e adiciona transicao:

\begin{lstlisting}[style=pascal]
SetLength(localDFATransitions, Length(...) + 1);
localDFATransitions[High(...)].FromState := curSet;
localDFATransitions[High(...)].Symbol := sym;
localDFATransitions[High(...)].ToState := nextSet;
\end{lstlisting}

Armazena em DOIS lugares:
\begin{itemize}
    \item \texttt{localDFATransitions} (variavel local)
    \item \texttt{DFATransitions} (campo da classe)
\end{itemize}

\subsection{Descoberta de Novos Estados (L3426-3439)}

\textbf{Referencia:} \texttt{MainForm.pas} linhas 3426-3439

\textbf{Condicao:} Estado e novo SE:
\begin{itemize}
    \item \texttt{dfaMap.IndexOf(key) = -1} (nao esta no mapa)
    \item \texttt{key $\neq$ '\{\}' } (nao e estado de erro)
\end{itemize}

\textbf{Acoes:}
\begin{lstlisting}[style=pascal]
cloneSet := nextSet.Clone;
dfaMap.AddObject(key, cloneSet);
localDFAStates.Add(key);
workQ.Add(key);  (* Adicionar a fila! *)
\end{lstlisting}

\textbf{Por que Clone?} \texttt{nextSet} sera liberado no \texttt{finally}. Precisamos de copia!

\subsubsection{Gerenciamento de Memoria Detalhado}

\textbf{Problema sem Clone():}
\begin{lstlisting}[style=pascal]
nextSet := TStringSet.Create;  (* Cria nextSet *)
dfaMap.AddObject(key, nextSet);  (* Adiciona ao mapa *)
nextSet.Free;  (* ERRO! dfaMap aponta para memoria liberada! *)
\end{lstlisting}

\textbf{Por que nextSet precisa ser liberado?}
\begin{itemize}
    \item Criado dentro do loop: \texttt{nextSet := TStringSet.Create}
    \item Sem Free(): memory leak em CADA iteracao!
    \item AFD com 100 estados = 100 conjuntos nao liberados
\end{itemize}

\textbf{Solucao: Clone antes de adicionar}
\begin{lstlisting}[style=pascal]
nextSet := TStringSet.Create;     (* Conjunto temporario *)
cloneSet := nextSet.Clone;        (* Copia independente *)
dfaMap.AddObject(key, cloneSet);  (* Adiciona copia ao mapa *)
nextSet.Free;                     (* Libera temporario *)
(* cloneSet sera liberado quando dfaMap for destruido *)
\end{lstlisting}

\textbf{Ciclo de vida dos objetos:}
\begin{enumerate}
    \item \texttt{nextSet}: Vive apenas durante uma iteracao do loop
    \item \texttt{cloneSet}: Vive ate \texttt{dfaMap.Free}
    \item \texttt{dfaMap}: Liberado no bloco \texttt{finally} ao fim da funcao
\end{enumerate}

\subsubsection{Por Que IndexOf ao Inves de Contains?}

\textbf{TStringList.IndexOf vs busca manual:}
\begin{itemize}
    \item \texttt{IndexOf(key)} retorna indice (-1 se nao existe)
    \item Com \texttt{Sorted := True}, usa busca binaria $O(\log n)$
    \item Loop manual seria $O(n)$
    \item Com 1000 estados: $\log_2(1000) \approx 10$ vs 1000 comparacoes!
\end{itemize}

\textbf{Por que comparar com -1?}
\begin{itemize}
    \item IndexOf retorna -1 quando nao encontra
    \item Indices validos: 0, 1, 2, ..., Count-1
    \item \texttt{IndexOf(key) = -1} $\equiv$ "key nao existe"
\end{itemize}

\subsection{Formatacao da Saida (L3549-3638)}

\textbf{Referencia:} \texttt{MainForm.pas} linhas 3549-3638

Apresenta resultado no \texttt{memoOutput} com:
\begin{itemize}
    \item Cabecalho decorativo
    \item Alfabeto
    \item Lista de estados do AFD
    \item Estado inicial (sempre o primeiro)
    \item Estados finais
    \item Lista de transicoes
\end{itemize}

\textbf{Tambem copia para campos da classe:} DFAStates, DFAInitial, DFAFinals.

\subsubsection{Por Que Dois Destinos para o Resultado?}

\textbf{Destino 1 - memoOutput (UI):}
\begin{itemize}
    \item Usuario ve resultado imediatamente
    \item Formato texto legivel
    \item Pode copiar/colar
    \item Pode salvar em arquivo
\end{itemize}

\textbf{Destino 2 - Campos da classe (DFAStates, etc):}
\begin{itemize}
    \item Proxima etapa (minimizacao) precisa dos dados
    \item Funcao DrawAutomaton precisa desenhar diagrama
    \item Botao "Usar como entrada" precisa copiar
\end{itemize}

\textbf{Exemplo de pipeline:}
\begin{enumerate}
    \item Conversao AFN $\rightarrow$ AFD: resultado em DFAStates
    \item Usuario clica "Minimizar": le DFAStates, escreve MinDFAStates
    \item Usuario clica "Desenhar": le DFAStates, desenha diagrama
    \item Usuario clica "Exportar": le memoOutput, salva arquivo
\end{enumerate}

\textbf{Por que nao armazenar APENAS em UI?}
\begin{itemize}
    \item Parse de texto e custoso ($O(n \cdot m)$)
    \item Parsing pode introduzir erros
    \item Estruturas de dados ja existem, reaproveitar!
\end{itemize}

\subsubsection{Formato de Estados do AFD}

\textbf{Notacao escolhida:} "\{q0,q1,q2\}"

\textbf{Por que chaves?}
\begin{itemize}
    \item Notacao matematica padrao para conjuntos
    \item Usuario reconhece: "isso e um conjunto!"
    \item Diferencia de estados simples: "q0" vs "\{q0\}"
\end{itemize}

\textbf{Por que virgulas sem espacos?}
\begin{itemize}
    \item ToString() de TStringSet usa CommaText
    \item CommaText automaticamente coloca virgulas
    \item Formato compacto: "\{q0,q1\}" vs "\{q0, q1\}"
\end{itemize}

\textbf{Alternativas consideradas:}
\begin{itemize}
    \item "(q0,q1)" -- parece tupla, nao conjunto
    \item "[q0,q1]" -- parece array
    \item "q0q1" -- nao distingue "q01" de "q0,q1"
\end{itemize}

\subsection{Habilitacao da Interface (L3640-3647)}

\textbf{Referencia:} \texttt{MainForm.pas} linhas 3640-3647

\begin{lstlisting}[style=pascal]
btnMinimize.Enabled := True;
PageControl2.ActivePage := TabOutput;
PaintBoxNFA.Invalidate;
PaintBoxDFA.Invalidate;
\end{lstlisting}

\textbf{Acoes:}
\begin{itemize}
    \item Habilita botao "Minimizar AFD"
    \item Muda para aba de resultado
    \item Redesenha diagramas
\end{itemize}

\subsection{Limpeza de Memoria (L3650-3665)}

\textbf{Referencia:} \texttt{MainForm.pas} linhas 3650-3665

Libera TODAS estruturas locais no \texttt{finally}:

\begin{lstlisting}[style=pascal]
localDFAStates.Free;
workQ.Free;
dfaMap.Free;  (* Libera TStringSets tambem! *)
isFinal.Free;
Alphabet.Free;
States.Free;
Initials.Free;
Finals.Free;
parts.Free;
\end{lstlisting}

\textbf{IMPORTANTE:} \texttt{dfaMap.OwnsObjects = True} $\rightarrow$ Free() tambem libera todos os TStringSet!

\begin{table}[ht]
\centering
\small
\begin{tabular}{|l|c|p{5cm}|}
\hline
\textbf{Funcao} & \textbf{Linhas} & \textbf{Descricao} \\
\hline
\multicolumn{3}{|c|}{\textbf{PARTE 1: Estruturas e Interface (L1-2254)}} \\
\hline
TTransition (type) & 39-43 & Record de transicao \\
TTransitionArray (type) & 49 & Array dinamico \\
Variaveis globais & 135-165 & AFN-e, AFN, AFD, MinDFA \\
TStringSet (class) & 299-311 & Declaracao da classe \\
TStringSet.Create & 333-337 & Construtor \\
TStringSet.Destroy & 339-343 & Destrutor \\
TStringSet.Add & 360-367 & Adicionar elemento \\
TStringSet.Contains & 377-381 & Verificar existencia \\
TStringSet.Clone & 398-410 & Clonar conjunto \\
TStringSet.IsEmpty & 414-418 & Verificar vazio \\
TStringSet.ToString & 431-445 & Converter para string \\
TStringSet.Count & 449-453 & Numero de elementos \\
TStringSet.Item & 464-475 & Acessar por indice \\
FormCreate & 488-571 & Inicializar interface \\
LoadTestFiles & 586-643 & Scan arquivos testes/ \\
OnTestFileSelected & 645-679 & Handler ComboBox \\
ConvertJSONToText & 693-812 & Converter arquivo JSON \\
ConvertJSONTextToText & 829-918 & Converter string JSON \\
LoadSampleFile & 952-1043 & Carregar exemplo \\
btnLoadFileClick & 1085-1170 & Handler carregar arquivo \\
btnClearClick & 1237-1360 & Limpar interface \\
btnRemoveEpsilonClick & 1395-1566 & Handler remocao epsilon \\
btnUseAsInputClick & 1603-1720 & Copiar AFN para input \\
btnConvertClick & 1863-2068 & Handler conversao AFN para AFD \\
\hline
\multicolumn{3}{|c|}{\textbf{PARTE 2: Algoritmo AFN para AFD (L2958-3795)}} \\
\hline
ConvertAFNtoAFD (inicio) & 2958-3020 & Declaracoes e variaveis \\
Limpeza e criacao & 2970-3010 & Liberar e criar estruturas \\
Validacao entrada & 3011-3018 & Verificar 4+ linhas \\
Parse alfabeto & 3040-3058 & Ler linha 0 (simbolos) \\
Parse estados & 3060-3068 & Ler linha 1 (estados) \\
Parse iniciais & 3070-3078 & Ler linha 2 (iniciais) \\
Parse finais & 3080-3088 & Ler linha 3 (finais) \\
Parse transicoes & 3090-3220 & Ler linhas 4+ (transicoes) \\
Logging entrada & 3222-3235 & Exibir AFN no console \\
Inicializacao BFS & 3237-3258 & Criar estruturas AFD \\
Estado inicial AFD & 3260-3270 & Conjunto dos iniciais \\
Loop principal BFS & 3272-3580 & Subset construction \\
Determinar finais & 3360-3380 & Marcar estados finais \\
Gerar transicoes & 3390-3500 & Para cada simbolo \\
Descobrir novos estados & 3510-3550 & Adicionar a fila \\
Formatacao resultado & 3582-3680 & Texto para memoOutput \\
Atualizar interface & 3682-3696 & Habilitar botoes \\
PaintBoxEpsilonNFAPaint & 3698-3723 & Desenhar AFN-epsilon \\
PaintBoxNFAPaint & 3725-3737 & Desenhar AFN \\
PaintBoxDFAPaint & 3739-3769 & Desenhar AFD \\
PaintBoxMinDFAPaint & 3771-3792 & Desenhar MinDFA \\
\hline
\end{tabular}
\caption{Indice Completo com Linhas Exatas}
\end{table}

\subsection{Paint Handlers - Desenho de Diagramas (L3698-3792)}

\textbf{Referencia:} \texttt{MainForm.pas} linhas 3698-3792

O codigo implementa 4 event handlers para desenhar diagramas de automatos.

\subsubsection{PaintBoxEpsilonNFAPaint (L3698-3710)}

\textbf{Proposito:} Desenhar AFN-$\varepsilon$ no painel grafico.

\textbf{Quando e chamado:}
\begin{itemize}
    \item Primeira exibicao da aba
    \item Apos \texttt{Invalidate()}
    \item Redimensionamento da janela
\end{itemize}

\begin{lstlisting}[style=pascal]
procedure TFormMain.PaintBoxEpsilonNFAPaint(Sender: TObject);
begin
  if not Assigned(EpsilonNFAStates) then Exit;
  DrawAutomaton(PaintBoxEpsilonNFA.Canvas, 
                EpsilonNFAStates, EpsilonNFAInitials,
                EpsilonNFAFinals, EpsilonNFATransitions);
end;
\end{lstlisting}

\textbf{Validacao:} Se nao ha dados, sai sem desenhar.

\subsubsection{PaintBoxNFAPaint (L3712-3718)}

\textbf{Proposito:} Desenhar AFN (sem epsilon) no painel.

Mesmo padrao do anterior, mas usa estruturas \texttt{NFAXXX}.

\subsubsection{PaintBoxDFAPaint (L3720-3741)}

\textbf{Proposito:} Desenhar AFD no painel.

\textbf{Diferenca importante:} AFD tem apenas 1 estado inicial!

\begin{lstlisting}[style=pascal]
procedure TFormMain.PaintBoxDFAPaint(Sender: TObject);
var Initials: TStringList;
begin
  if not Assigned(DFAStates) then Exit;
  
  Initials := TStringList.Create;
  try
    if DFAStates.Count > 0 then
      Initials.Add(DFAStates[0]);  (* Primeiro = inicial *)
    
    DrawAutomaton(PaintBoxDFA.Canvas, DFAStates, 
                  Initials, DFAFinals, DFATransitions);
  finally
    Initials.Free;  (* Liberar memoria! *)
  end;
end;
\end{lstlisting}

\textbf{Por que TStringList temporaria?} DrawAutomaton espera TStringList de iniciais. AFD tem string unica, precisa converter.

\subsubsection{PaintBoxMinDFAPaint (L3743-3764)}

\textbf{Proposito:} Desenhar AFD Minimizado no painel.

Mesmo padrao do PaintBoxDFAPaint, mas usa:
\begin{itemize}
    \item \texttt{MinDFAStates}
    \item \texttt{MinDFAInitial}
    \item \texttt{MinDFAFinals}
    \item \texttt{MinDFATransitions}
\end{itemize}

\subsection{Complexidade do Algoritmo Completo}

\textbf{Analise de cada parte:}

\begin{table}[h]
\centering
\small
\begin{tabular}{|l|c|}
\hline
\textbf{Operacao} & \textbf{Complexidade} \\
\hline
Parse entrada & $O(n + m)$ \\
Estado inicial & $O(n_{init})$ \\
Loop principal BFS & $O(2^n)$ iteracoes \\
Por cada estado & $O(|\Sigma| \cdot |S| \cdot m)$ \\
Verificacao final & $O(|S| \cdot |F|)$ \\
Calculo destinos & $O(|S| \cdot m)$ \\
\hline
\textbf{Total pior caso} & $\mathbf{O(2^n \cdot |\Sigma| \cdot n \cdot m)}$ \\
\hline
\end{tabular}
\caption{Complexidade do Subset Construction}
\end{table}

Onde:
\begin{itemize}
    \item $n$ = numero de estados do AFN
    \item $m$ = numero de transicoes do AFN
    \item $|\Sigma|$ = tamanho do alfabeto
    \item $|S|$ = tamanho medio dos conjuntos
    \item $|F|$ = numero de estados finais
\end{itemize}

\textbf{Explosao exponencial:} No pior caso, AFD pode ter $2^n$ estados!

\textbf{Exemplo:} AFN com 10 estados $\rightarrow$ AFD com ate 1024 estados.

\subsection{Mapa Completo de Funcoes}

\begin{table}[h]
\centering
\small
\begin{tabular}{|l|c|p{6cm}|}
\hline
\textbf{Funcao} & \textbf{Linhas} & \textbf{Descricao} \\
\hline
\multicolumn{3}{|c|}{\textbf{PARTE 1: Estruturas e Interface}} \\
\hline
TTransition & 39-43 & Record de transicao \\
TStringSet.Create & 333-337 & Construtor \\
TStringSet.Destroy & 339-343 & Destrutor \\
FormCreate & 488-571 & Inicializar interface \\
LoadTestFiles & 586-643 & Scan arquivos testes/ \\
ConvertJSONToText & 693-812 & Converter arquivo JSON \\
ConvertJSONTextToText & 829-918 & Converter string JSON \\
LoadSampleFile & 952-1043 & Carregar exemplo \\
btnLoadFileClick & 1085-1170 & Handler carregar arquivo \\
btnClearClick & 1237-1360 & Limpar interface \\
btnRemoveEpsilonClick & 1395-1566 & Handler remocao epsilon \\
btnUseAsInputClick & 1603-1720 & Copiar AFN para input \\
btnConvertClick & 1863-2068 & Handler conversao AFN para AFD \\
GetSetName & 2006-2058 & TStringSet para String \\
ComputeEpsilonClosure & 2109-2300 & Calcular epsilon-fecho \\
\hline
\multicolumn{3}{|c|}{\textbf{PARTE 2: Algoritmo AFN para AFD}} \\
\hline
ConvertAFNtoAFD (inicio) & 2958-3020 & Declaracoes e init \\
Parse alfabeto & 3040-3055 & Ler linha 0 \\
Parse estados & 3057-3065 & Ler linha 1 \\
Parse iniciais & 3067-3075 & Ler linha 2 \\
Parse finais & 3077-3085 & Ler linha 3 \\
Parse transicoes & 3150-3220 & Ler linhas 4+ \\
Inicializacao & 3221-3270 & Criar estruturas AFD \\
Loop principal & 3271-3600 & BFS subset construction \\
Formatacao saida & 3601-3720 & Gerar memoOutput \\
PaintBoxEpsilonNFAPaint & 3721-3735 & Desenhar AFN-epsilon \\
PaintBoxNFAPaint & 3737-3745 & Desenhar AFN \\
PaintBoxDFAPaint & 3747-3765 & Desenhar AFD \\
PaintBoxMinDFAPaint & 3767-3785 & Desenhar MinDFA \\
\hline
\end{tabular}
\caption{Indice Completo de Funcoes - Secoes do Henrique}
\end{table}

\newpage
\section{Conclusao}

\subsection{Resumo das Responsabilidades}

\textbf{Henrique implementou:}

\begin{enumerate}
    \item \textbf{Estruturas de Dados Fundamentais}
    \begin{itemize}
        \item TTransition (record de transicao)
        \item TStringSet (classe de conjuntos)
        \item Variaveis globais para AFN/AFD
    \end{itemize}
    
    \item \textbf{Interface e Carregamento}
    \begin{itemize}
        \item Inicializacao do formulario
        \item Scan de arquivos de teste
        \item Parser JSON para TXT
    \end{itemize}
    
    \item \textbf{Algoritmo Principal: AFN para AFD}
    \begin{itemize}
        \item Construcao de subconjuntos
        \item Determinacao de estados finais
        \item Geracao de transicoes deterministicas
        \item Logging e interface automatica
    \end{itemize}
\end{enumerate}

\subsection{Pontos-Chave para Apresentacao}

\begin{enumerate}
    \item \textbf{TStringSet e essencial} - evita duplicatas nos conjuntos
    \item \textbf{Cada estado do AFD = conjunto de estados do AFN}
    \item \textbf{Loop principal:} Para cada conjunto nao processado, para cada simbolo, calcule TODOS os destinos possiveis
    \item \textbf{Estado final:} Se o conjunto contem pelo menos um estado final do AFN
    \item \textbf{Complexidade:} Exponencial no pior caso ($2^n$), mas tipico e linear
    \item \textbf{Parse robusto:} Aceita JSON e TXT, converte automaticamente
    \item \textbf{Validacao completa:} Detecta epsilon e orienta usuario sobre workflow correto
\end{enumerate}

\subsection{Complexidade do Algoritmo}

\begin{table}[h]
\centering
\begin{tabular}{|l|c|l|}
\hline
\textbf{Operacao} & \textbf{Complexidade} & \textbf{Motivo} \\
\hline
Estados do AFD & $O(2^n)$ & No pior caso, todos subconjuntos \\
Loop principal & $O(estados \cdot |\Sigma|)$ & Para cada estado/simbolo \\
Busca de transicoes & $O(m)$ & $m$ = transicoes do AFN \\
\hline
\textbf{Total} & $\mathbf{O(2^n \cdot |\Sigma| \cdot m)}$ & Pior caso exponencial \\
\hline
\end{tabular}
\caption{Analise de Complexidade}
\end{table}

\end{document}
